%% The last subsection of the appendix "An Introduction to the Framework":
%% the end of Section 2 of the vqupit paper
%% (vqupit/sections/constructions.lagda.tex).  make.sh writes
%% ../latex/pldi/lagda/app-constructions.tex; edit this file.
\begin{code}[hide]
module pldi.lagda.app-constructions where

open import Data.Sum using (_⊎_)
open import Level using (0ℓ)
open import Algebra.Bundles using (Group)
open import Word.Base using (Word ; WRel ; _ʰ' ; _ⁿ')
open import Presentation.Definitions using (_IsPresentationOf_)
open import Presentation.Construct.Base using ([_]ₗ ; [_]ᵣ ; ConjRelʷ ; EmptyRel)

infix 4 _≈₁_ _===₁_ _===₂_
infix 5 _⋄_⋄_

postulate
  proof-elided : ∀ {ℓ} {A : Set ℓ} → A
  A B N H : Set
  u v : Word A
  _===₁_ : WRel N
  _===₂_ : WRel H
  _≈₁_ : Word N → Word N → Set
  Γ Δ : Set
  conj : H → N → Word N
  G1⋊G2 : Group 0ℓ 0ℓ
  RelTwist : Set → Set → Set
  _∪_ : Set → Set → Set
\end{code}
\subsection{Products of Presentations, and Their Presentation Theorems}\label{app:constructions}

Presentations such as the one just obtained can be combined.  Relation
sets over a disjoint union of alphabets are built from a single
primitive, a three-way join whose \emph{mixed} component is the design
parameter:

\begin{code}
data _⋄_⋄_ (Γ₁ : WRel A) (Γ₂ : WRel B) (Γ₃ : WRel (A ⊎ B)) : WRel (A ⊎ B) where
  left  : Γ₁ u v → (Γ₁ ⋄ Γ₂ ⋄ Γ₃) [ u ]ₗ [ v ]ₗ
  right : Γ₂ u v → (Γ₁ ⋄ Γ₂ ⋄ Γ₃) [ u ]ᵣ [ v ]ᵣ
  mid   : Γ₃ u v → (Γ₁ ⋄ Γ₂ ⋄ Γ₃) u v
\end{code}

Choosing the mixed relation recovers the classical constructions: free
product ($\aid{EmptyRel}$), direct product ($\aid{CommRel}$), semidirect
product ($\aid{ConjRelʷ}\;\mathit{conj}$: moving a letter $h$ of the
acting group past a letter $n$ of the normal subgroup conjugates $n$, by
a \emph{word}, $h \bullet n = \mathit{conj}\;h\;n \bullet h$) and
definitional extension ($\aid{SugarRel}$).  For each, the library lifts
normal forms from the factors to the construction and proves that if
the factors present concrete groups, the constructed relation set
presents the constructed group.  Beyond presentations of the factors,
the semidirect theorem needs two well-definedness conditions on the
action: that it respects
the acting group's axioms in its first argument (\aid{conj-hyph}) and
the normal subgroup's in its second (\aid{conj-hypn}):

\begingroup\renewcommand{\AgdaCodeStyle}{\footnotesize}
\begin{code}
module _ (conj-hyph : ∀ {c d} n → c ===₂ d → (conj ʰ') c n ≈₁ (conj ʰ') d n)
         (conj-hypn : ∀ c {w v} → w ===₁ v → (conj ⁿ') c w ≈₁ (conj ⁿ') c v) where
  module Presentation (G1 G2 : Group 0ℓ 0ℓ)
    (p1 : Γ IsPresentationOf G1) (p2 : Δ IsPresentationOf G2) where
    dpres : (Γ ⋄ Δ ⋄ ConjRelʷ conj) IsPresentationOf G1⋊G2
\end{code}
\endgroup
\begin{code}[hide]
    dpres = proof-elided

extension-presentation : Set → Set → Set → Set → Set
\end{code}

\noindent Here \aid{G1⋊G2} is the semidirect product of the two
\emph{semantic} groups, built from an \aid{Action} record whose action
is the syntactic \aid{conj} transported through the two presentation
isomorphisms.  A further module proves the classical recipe for
presenting a group \emph{extension} $1 \to N \to G \to Q \to 1$ from
presentations of $N$ and $Q$, how $G$ conjugates $N$, and the
correction in $N$ by which each lifted relator of $Q$ fails to hold
(cf.\ Proposition~2.55 of~\cite{holt2005handbook}): a second mixed
component, \aid{RelTwist}, twists each relator of the quotient by its
correction, and the relation set of the recipe is
\begingroup\renewcommand{\AgdaCodeStyle}{\footnotesize}
\begin{code}
extension-presentation S R conj corr = S ⋄ EmptyRel ⋄ (ConjRelʷ conj ∪ RelTwist R corr)
\end{code}
\endgroup
Semidirect products are the special case with trivial corrections.
With these theorems a presentation of a product or of an extension is
\emph{assembled} from presentations of its factors, with no fresh coset
enumeration at all.
